% Shared preamble for Wolf lecture note transcriptions
% Usage: % Shared preamble for Wolf lecture note transcriptions
% Usage: % Shared preamble for Wolf lecture note transcriptions
% Usage: % Shared preamble for Wolf lecture note transcriptions
% Usage: \input{preamble} in each chapter file
% Include guard: safe to \input multiple times (e.g. from main_wolf.tex)
\ifdefined\wolfpreambleloaded\endinput\fi
\def\wolfpreambleloaded{}

\usepackage{amsmath,amssymb,amsthm}
% Match Wolf book-class layout: a4paper, ~345pt text width
\usepackage[a4paper, textwidth=345pt, textheight=550pt]{geometry}
\usepackage{hyperref}
\usepackage{mathtools}

% ── Notation (consistent across all chapters) ──
\newcommand{\mcM}{\mathcal{M}}       % matrix algebra M_d(C)
\newcommand{\mcB}{\mathcal{B}}       % bounded operators B(H)
\newcommand{\mcA}{\mathcal{A}}       % multiplicative domain
\newcommand{\mcH}{\mathcal{H}}       % Hilbert space H
\newcommand{\diag}{\operatorname{diag}}
\newcommand{\one}{\mathbb{1}}        % identity operator
\newcommand{\CC}{\mathbb{C}}
\newcommand{\RR}{\mathbb{R}}
\newcommand{\NN}{\mathbb{N}}
\newcommand{\ZZ}{\mathbb{Z}}
\newcommand{\tr}{\operatorname{tr}}
\newcommand{\spec}{\operatorname{spec}}
\newcommand{\rank}{\operatorname{rank}}
\newcommand{\supp}{\operatorname{supp}}
\newcommand{\range}{\operatorname{range}}
\newcommand{\spn}{\operatorname{span}}
\newcommand{\FT}{\mathcal{F}_T}      % fixed point set
\newcommand{\XT}{\mathcal{X}_T}      % peripheral eigenspace
\newcommand{\mcX}{\mathcal{X}}       % peripheral eigenspace (bare)
\newcommand{\id}{\mathrm{id}}        % identity map
\newcommand{\ket}[1]{|#1\rangle}
\newcommand{\bra}[1]{\langle#1|}
\newcommand{\braket}[2]{\langle#1|#2\rangle}
\newcommand{\ketbra}[2]{|#1\rangle\!\langle#2|}

% ── Helper: properly undefine a theorem environment for redefinition ──
% Usage: \UndefEnv{theorem} before \newtheorem{theorem}{...}
\makeatletter
\newcommand{\UndefEnv}[1]{%
  % Remove from section reset list (if registered there)
  \@removefromreset{#1}{section}%
  \expandafter\let\csname #1\endcsname\relax
  \expandafter\let\csname end#1\endcsname\relax
  \expandafter\let\csname c@#1\endcsname\relax
}
\makeatother

% ── Theorem environments ──
\theoremstyle{plain}
\newtheorem{theorem}{Theorem}[section]
\newtheorem{proposition}[theorem]{Proposition}
\newtheorem{corollary}[theorem]{Corollary}
\newtheorem{lemma}[theorem]{Lemma}
\theoremstyle{definition}
\newtheorem{definition}[theorem]{Definition}
\newtheorem{example}[theorem]{Example}
\newtheorem{problem}[theorem]{Problem}
\theoremstyle{remark}
\newtheorem{remark}[theorem]{Remark} in each chapter file
% Include guard: safe to \input multiple times (e.g. from main_wolf.tex)
\ifdefined\wolfpreambleloaded\endinput\fi
\def\wolfpreambleloaded{}

\usepackage{amsmath,amssymb,amsthm}
% Match Wolf book-class layout: a4paper, ~345pt text width
\usepackage[a4paper, textwidth=345pt, textheight=550pt]{geometry}
\usepackage{hyperref}
\usepackage{mathtools}

% ── Notation (consistent across all chapters) ──
\newcommand{\mcM}{\mathcal{M}}       % matrix algebra M_d(C)
\newcommand{\mcB}{\mathcal{B}}       % bounded operators B(H)
\newcommand{\mcA}{\mathcal{A}}       % multiplicative domain
\newcommand{\mcH}{\mathcal{H}}       % Hilbert space H
\newcommand{\diag}{\operatorname{diag}}
\newcommand{\one}{\mathbb{1}}        % identity operator
\newcommand{\CC}{\mathbb{C}}
\newcommand{\RR}{\mathbb{R}}
\newcommand{\NN}{\mathbb{N}}
\newcommand{\ZZ}{\mathbb{Z}}
\newcommand{\tr}{\operatorname{tr}}
\newcommand{\spec}{\operatorname{spec}}
\newcommand{\rank}{\operatorname{rank}}
\newcommand{\supp}{\operatorname{supp}}
\newcommand{\range}{\operatorname{range}}
\newcommand{\spn}{\operatorname{span}}
\newcommand{\FT}{\mathcal{F}_T}      % fixed point set
\newcommand{\XT}{\mathcal{X}_T}      % peripheral eigenspace
\newcommand{\mcX}{\mathcal{X}}       % peripheral eigenspace (bare)
\newcommand{\id}{\mathrm{id}}        % identity map
\newcommand{\ket}[1]{|#1\rangle}
\newcommand{\bra}[1]{\langle#1|}
\newcommand{\braket}[2]{\langle#1|#2\rangle}
\newcommand{\ketbra}[2]{|#1\rangle\!\langle#2|}

% ── Helper: properly undefine a theorem environment for redefinition ──
% Usage: \UndefEnv{theorem} before \newtheorem{theorem}{...}
\makeatletter
\newcommand{\UndefEnv}[1]{%
  % Remove from section reset list (if registered there)
  \@removefromreset{#1}{section}%
  \expandafter\let\csname #1\endcsname\relax
  \expandafter\let\csname end#1\endcsname\relax
  \expandafter\let\csname c@#1\endcsname\relax
}
\makeatother

% ── Theorem environments ──
\theoremstyle{plain}
\newtheorem{theorem}{Theorem}[section]
\newtheorem{proposition}[theorem]{Proposition}
\newtheorem{corollary}[theorem]{Corollary}
\newtheorem{lemma}[theorem]{Lemma}
\theoremstyle{definition}
\newtheorem{definition}[theorem]{Definition}
\newtheorem{example}[theorem]{Example}
\newtheorem{problem}[theorem]{Problem}
\theoremstyle{remark}
\newtheorem{remark}[theorem]{Remark} in each chapter file
% Include guard: safe to \input multiple times (e.g. from main_wolf.tex)
\ifdefined\wolfpreambleloaded\endinput\fi
\def\wolfpreambleloaded{}

\usepackage{amsmath,amssymb,amsthm}
% Match Wolf book-class layout: a4paper, ~345pt text width
\usepackage[a4paper, textwidth=345pt, textheight=550pt]{geometry}
\usepackage{hyperref}
\usepackage{mathtools}

% ── Notation (consistent across all chapters) ──
\newcommand{\mcM}{\mathcal{M}}       % matrix algebra M_d(C)
\newcommand{\mcB}{\mathcal{B}}       % bounded operators B(H)
\newcommand{\mcA}{\mathcal{A}}       % multiplicative domain
\newcommand{\mcH}{\mathcal{H}}       % Hilbert space H
\newcommand{\diag}{\operatorname{diag}}
\newcommand{\one}{\mathbb{1}}        % identity operator
\newcommand{\CC}{\mathbb{C}}
\newcommand{\RR}{\mathbb{R}}
\newcommand{\NN}{\mathbb{N}}
\newcommand{\ZZ}{\mathbb{Z}}
\newcommand{\tr}{\operatorname{tr}}
\newcommand{\spec}{\operatorname{spec}}
\newcommand{\rank}{\operatorname{rank}}
\newcommand{\supp}{\operatorname{supp}}
\newcommand{\range}{\operatorname{range}}
\newcommand{\spn}{\operatorname{span}}
\newcommand{\FT}{\mathcal{F}_T}      % fixed point set
\newcommand{\XT}{\mathcal{X}_T}      % peripheral eigenspace
\newcommand{\mcX}{\mathcal{X}}       % peripheral eigenspace (bare)
\newcommand{\id}{\mathrm{id}}        % identity map
\newcommand{\ket}[1]{|#1\rangle}
\newcommand{\bra}[1]{\langle#1|}
\newcommand{\braket}[2]{\langle#1|#2\rangle}
\newcommand{\ketbra}[2]{|#1\rangle\!\langle#2|}

% ── Helper: properly undefine a theorem environment for redefinition ──
% Usage: \UndefEnv{theorem} before \newtheorem{theorem}{...}
\makeatletter
\newcommand{\UndefEnv}[1]{%
  % Remove from section reset list (if registered there)
  \@removefromreset{#1}{section}%
  \expandafter\let\csname #1\endcsname\relax
  \expandafter\let\csname end#1\endcsname\relax
  \expandafter\let\csname c@#1\endcsname\relax
}
\makeatother

% ── Theorem environments ──
\theoremstyle{plain}
\newtheorem{theorem}{Theorem}[section]
\newtheorem{proposition}[theorem]{Proposition}
\newtheorem{corollary}[theorem]{Corollary}
\newtheorem{lemma}[theorem]{Lemma}
\theoremstyle{definition}
\newtheorem{definition}[theorem]{Definition}
\newtheorem{example}[theorem]{Example}
\newtheorem{problem}[theorem]{Problem}
\theoremstyle{remark}
\newtheorem{remark}[theorem]{Remark} in each chapter file
% Include guard: safe to \input multiple times (e.g. from main_wolf.tex)
\ifdefined\wolfpreambleloaded\endinput\fi
\def\wolfpreambleloaded{}

\usepackage{amsmath,amssymb,amsthm}
% Match Wolf book-class layout: a4paper, ~345pt text width
\usepackage[a4paper, textwidth=345pt, textheight=550pt]{geometry}
\usepackage{hyperref}
\usepackage{mathtools}

% ── Notation (consistent across all chapters) ──
\newcommand{\mcM}{\mathcal{M}}       % matrix algebra M_d(C)
\newcommand{\mcB}{\mathcal{B}}       % bounded operators B(H)
\newcommand{\mcA}{\mathcal{A}}       % multiplicative domain
\newcommand{\mcH}{\mathcal{H}}       % Hilbert space H
\newcommand{\diag}{\operatorname{diag}}
\newcommand{\one}{\mathbb{1}}        % identity operator
\newcommand{\CC}{\mathbb{C}}
\newcommand{\RR}{\mathbb{R}}
\newcommand{\NN}{\mathbb{N}}
\newcommand{\ZZ}{\mathbb{Z}}
\newcommand{\tr}{\operatorname{tr}}
\newcommand{\spec}{\operatorname{spec}}
\newcommand{\rank}{\operatorname{rank}}
\newcommand{\supp}{\operatorname{supp}}
\newcommand{\range}{\operatorname{range}}
\newcommand{\spn}{\operatorname{span}}
\newcommand{\FT}{\mathcal{F}_T}      % fixed point set
\newcommand{\XT}{\mathcal{X}_T}      % peripheral eigenspace
\newcommand{\mcX}{\mathcal{X}}       % peripheral eigenspace (bare)
\newcommand{\id}{\mathrm{id}}        % identity map
\newcommand{\ket}[1]{|#1\rangle}
\newcommand{\bra}[1]{\langle#1|}
\newcommand{\braket}[2]{\langle#1|#2\rangle}
\newcommand{\ketbra}[2]{|#1\rangle\!\langle#2|}

% ── Helper: properly undefine a theorem environment for redefinition ──
% Usage: \UndefEnv{theorem} before \newtheorem{theorem}{...}
\makeatletter
\newcommand{\UndefEnv}[1]{%
  % Remove from section reset list (if registered there)
  \@removefromreset{#1}{section}%
  \expandafter\let\csname #1\endcsname\relax
  \expandafter\let\csname end#1\endcsname\relax
  \expandafter\let\csname c@#1\endcsname\relax
}
\makeatother

% ── Theorem environments ──
\theoremstyle{plain}
\newtheorem{theorem}{Theorem}[section]
\newtheorem{proposition}[theorem]{Proposition}
\newtheorem{corollary}[theorem]{Corollary}
\newtheorem{lemma}[theorem]{Lemma}
\theoremstyle{definition}
\newtheorem{definition}[theorem]{Definition}
\newtheorem{example}[theorem]{Example}
\newtheorem{problem}[theorem]{Problem}
\theoremstyle{remark}
\newtheorem{remark}[theorem]{Remark}